% GENERATED FILE. Edit canonical lessons, then run npm run generate:books.
% canonical-source-hash: fnv1a64:1b7ec0c8838c2bcb
% canonical-lessons: ZH-C294-qingkuang, ZH-C294-yinqi, ZH-C294-yingxiang, ZH-C294-zhongyu, ZH-C294-yushi

\chapter{How It Turned Out}
\label{ch:zh-how-it-turned-out}
\hlchaptermodality{294}

\begin{chapteropening}
\textbf{By the end of this chapter:} \emph{I can say how things stood, what caused something, what it changed and how it ended.}

Complete the last lesson of chapter 294: I can say how things stood, what caused something, what it changed and how it ended.
\end{chapteropening}

\section[qíngkuàng]{qíngkuàng — situation; how things stand}

\label{lesson:ZH-C294-qingkuang}



\begin{quote}
Before the new one: say the Mandarin for the front-page story, then the Mandarin for to follow a story closely.
\end{quote}

\subsection*{You'll want to know: qíngkuàng}
\emph{qíngkuàng} — \textquotedblleft{}situation; how things stand\textquotedblright{}.

\begin{quote}
\emph{nǐ děi gàosu wǒ qíngkuàng, bùrán wǒ huì dānxīn} — You have to tell me how things are, or I'll worry.
\end{quote}

Its characters come later; for now it is a word for the ear.

\begin{cousinweb}
\emph{qíngkuàng} is built from \emph{qíng} 'feeling' and \emph{kuàng} 'to compare', both inherited from Middle Chinese, \emph{dzjeng} and \emph{xjwangH} (Guangyun, in Baxter's transcription), which glosses the first 'what a person desires' and the second 'to compare'; how the pair came to mean 'a situation' is not settled. \emph{xjwangH} would regularly give huàng; the k of kuàng is irregular.
\end{cousinweb}

\begin{grammarlens}[title={What you've built}]
One more word for telling what happened.
\end{grammarlens}

\subsection*{Your turn}
\begin{itemize}
\raggedright
  \item \emph{Say it:} \emph{qíngkuàng}
  \item \emph{Say it:} \emph{qíngkuàng}, once more
  \item \emph{Say it:} \emph{nǐ děi gàosu wǒ qíngkuàng, bùrán wǒ huì dānxīn} — You have to tell me how things are, or I'll worry.
  \item \emph{Recall:} say the Mandarin for the front-page story, then the Mandarin for to follow a story closely, then say \emph{qíngkuàng} again
\end{itemize}

\begin{tcolorbox}[breakable,colback=teal!4,colframe=teal!35!black,title={Before you move on}]
What does qíngkuàng mean? (\textquotedblleft{}Situation; how things stand\textquotedblright{}.) Say it once more.
\end{tcolorbox}
\section[yǐnqǐ]{yǐnqǐ — to give rise to; to bring about}

\label{lesson:ZH-C294-yinqi}



\begin{quote}
Before the new one: say the Mandarin for to follow a story closely, then the Mandarin for how things stand.
\end{quote}

\subsection*{You'll want to know: yǐnqǐ}
\emph{yǐnqǐ} — \textquotedblleft{}to give rise to; to bring about\textquotedblright{}.

\begin{quote}
\emph{zhège xīnwén yǐnqǐ le dàjiā de zhùyì} — The story caught everyone's attention.
\end{quote}

Its characters come later; for now it is a word for the ear.

Two third tones meet in this word: yǐn is said with a rising tone before qǐ, and the spelling keeps both marks.

\begin{cousinweb}
\emph{yǐnqǐ} is built from \emph{yǐn} 'to draw (a bow)' and \zh{起} 'to rise', both inherited from Middle Chinese, \emph{yinX} and \emph{khiX} (Guangyun, in Baxter's transcription), which glosses the first 'long' and 'to draw a bow' and the second 'to rise' and 'to start'.
\end{cousinweb}

\begin{grammarlens}[title={What you've built}]
One more word for telling what happened.
\end{grammarlens}

\subsection*{Your turn}
\begin{itemize}
\raggedright
  \item \emph{Say it:} \emph{yǐnqǐ}
  \item \emph{Say it:} \emph{yǐnqǐ}, once more
  \item \emph{Say it:} \emph{zhège xīnwén yǐnqǐ le dàjiā de zhùyì} — The story caught everyone's attention.
  \item \emph{Recall:} say the Mandarin for to follow a story closely, then the Mandarin for how things stand, then say \emph{yǐnqǐ} again
\end{itemize}

\begin{tcolorbox}[breakable,colback=teal!4,colframe=teal!35!black,title={Before you move on}]
What does yǐnqǐ mean? (\textquotedblleft{}To give rise to; to bring about\textquotedblright{}.) Say it once more.
\end{tcolorbox}
\section[yǐngxiǎng]{yǐngxiǎng — to affect; influence; to have an effect on}

\label{lesson:ZH-C294-yingxiang}



\begin{quote}
Before the new one: say the Mandarin for how things stand, then the Mandarin for to give rise to something.
\end{quote}

\subsection*{You'll want to know: yǐngxiǎng}
\emph{yǐngxiǎng} — \textquotedblleft{}to affect; influence; to have an effect on\textquotedblright{}.

\begin{quote}
\emph{xiàyǔ yǐngxiǎng le bǐsài} — The rain affected the match.
\end{quote}

Its characters come later; for now it is a word for the ear.

Two third tones meet in this word: yǐng is said with a rising tone before xiǎng, and the spelling keeps both marks.

\begin{cousinweb}
\emph{yǐngxiǎng} is built from \emph{yǐng} 'a shadow' and \emph{xiǎng} 'a sound', both inherited from Middle Chinese, \emph{'jaengX} and \emph{xjangX} (Guangyun, in Baxter's transcription), which glosses the first 'a shape's shadow' and the second 'a sound'. It is usually explained as what follows a thing the way a shadow follows a shape and an echo a sound.
\end{cousinweb}

\begin{grammarlens}[title={What you've built}]
One more word for telling what happened.
\end{grammarlens}

\subsection*{Your turn}
\begin{itemize}
\raggedright
  \item \emph{Say it:} \emph{yǐngxiǎng}
  \item \emph{Say it:} \emph{yǐngxiǎng}, once more
  \item \emph{Say it:} \emph{xiàyǔ yǐngxiǎng le bǐsài} — The rain affected the match.
  \item \emph{Recall:} say the Mandarin for how things stand, then the Mandarin for to give rise to something, then say \emph{yǐngxiǎng} again
\end{itemize}

\begin{tcolorbox}[breakable,colback=teal!4,colframe=teal!35!black,title={Before you move on}]
What does yǐngxiǎng mean? (\textquotedblleft{}To affect; influence; to have an effect on\textquotedblright{}.) Say it once more.
\end{tcolorbox}
\section[zhōngyú]{zhōngyú — at long last}

\label{lesson:ZH-C294-zhongyu}



\begin{quote}
Before the new one: say the Mandarin for to give rise to something, then the Mandarin for to have an effect on.
\end{quote}

\subsection*{You'll want to know: zhōngyú}
\emph{zhōngyú} — \textquotedblleft{}at long last\textquotedblright{}.

\begin{quote}
\emph{děng le liǎng ge xiǎoshí, chē zhōngyú lái le} — After a two-hour wait, the bus came at last.
\end{quote}

Its characters come later; for now it is a word for the ear.

\begin{cousinweb}
\emph{zhōngyú} is built from \emph{zhōng} 'the end' and \emph{yú} 'at', both inherited from Middle Chinese, \emph{tsyuwng} and \emph{'jo} (Guangyun, in Baxter's transcription), which glosses the first 'the end' and 'to use up' and the second 'to dwell' and calls it a particle; its sense 'at' is not given there. The rising tone of yú is irregular for \emph{'jo}.
\end{cousinweb}

\begin{grammarlens}[title={What you've built}]
One more word for telling what happened.
\end{grammarlens}

\subsection*{Your turn}
\begin{itemize}
\raggedright
  \item \emph{Say it:} \emph{zhōngyú}
  \item \emph{Say it:} \emph{zhōngyú}, once more
  \item \emph{Say it:} \emph{děng le liǎng ge xiǎoshí, chē zhōngyú lái le} — After a two-hour wait, the bus came at last.
  \item \emph{Recall:} say the Mandarin for to give rise to something, then the Mandarin for to have an effect on, then say \emph{zhōngyú} again
\end{itemize}

\begin{tcolorbox}[breakable,colback=teal!4,colframe=teal!35!black,title={Before you move on}]
What does zhōngyú mean? (\textquotedblleft{}At long last\textquotedblright{}.) Say it once more.
\end{tcolorbox}
\section[yúshì]{yúshì — so; and then (as a result)}

\label{lesson:ZH-C294-yushi}



\begin{quote}
Before the new one: say the Mandarin for to have an effect on, then the Mandarin for at long last.
\end{quote}

\subsection*{You'll want to know: yúshì}
\emph{yúshì} — \textquotedblleft{}so; and then (as a result)\textquotedblright{}.

\begin{quote}
\emph{yǔ hěn dà, yúshì wǒmen huí jiā le} — It was raining hard, so we went home.
\end{quote}

Its characters come later; for now it is a word for the ear.

\begin{cousinweb}
\emph{yúshì} is built from \emph{yú} 'at' and \zh{是}, both inherited from Middle Chinese, \emph{'jo} and \emph{dzyeX} (Guangyun, in Baxter's transcription), which calls the first a particle, without the sense 'at', and glosses the second 'right, not wrong'; how the pair came to mean 'and so' is not settled. The rising tone of yú is irregular for \emph{'jo}.
\end{cousinweb}

\begin{grammarlens}[title={What you've built}]
One more word for telling what happened.
\end{grammarlens}

\subsection*{Your turn}
\begin{itemize}
\raggedright
  \item \emph{Say it:} \emph{yúshì}
  \item \emph{Say it:} \emph{yúshì}, once more
  \item \emph{Say it:} \emph{yǔ hěn dà, yúshì wǒmen huí jiā le} — It was raining hard, so we went home.
  \item \emph{Recall:} say the Mandarin for to have an effect on, then the Mandarin for at long last, then say \emph{yúshì} again
\end{itemize}

\begin{tcolorbox}[breakable,colback=teal!4,colframe=teal!35!black,title={Before you move on}]
What does yúshì mean? (\textquotedblleft{}So; and then (as a result)\textquotedblright{}.) Say it once more.
\end{tcolorbox}
